\subsection*{Partie A}

\begin{enumerate}
	\item La France a obtenu 11 médailles en 2010.

	\item
	\[\overline{m}=\dfrac{11+9+11+15+15+14}{6}=\dfrac{75}{6}=12,5.\]
	En moyenne, la France a obtenu 12,5 médailles lors de ces 6 jeux olympiques.

	\item
	\[\dfrac{13}{25}=0,52=\dfrac{52}{100}.\]
	En 2002, 52\,\% des médailles de la Norvège étaient des médailles d'or.
\end{enumerate}

\subsection*{Partie B}

\begin{enumerate}
	\item Les points A et H étant à la même altitude, H est à l'altitude \np{1814}~m donc :
	\[\text{HB}=\np{2359}-\np{1814}=545~\text{m}.\]

La différence d'altitude entre les points H et B est égale à $545$ m.

	\item Le triangle ABH est rectangle en H. D'après la trigonométrie :
	
\[\sin\left(\widehat{\text{BAH}}\right)=\dfrac{\text{HB}}{\text{AB}}=\dfrac{545}{\np{1600}}\]

donc $\widehat{\text{BAH}}\approx 19,9^{\circ}$ à la calculatrice.

Comme $19,9^{\circ}<30^{\circ}$, la norme de sécurité est respectée.

	\item

\[t=\dfrac{d}{v}=\dfrac{\np{1600}}{5}=320~\text{s}.\]

Or $320=5\times 60 + 20$, donc le trajet dure 5~min~20~s.
\end{enumerate}
